% --- Archon physics formalization source begin ---
% archon:physics
% archon:covers PhyXMiniProblems/problem_phyx_mini_0715.lean
% archon:source-report reports/phyx_mini/problem_phyx_mini_0715.source.json

\chapter{Physics problem phyx\_mini\_0715}
\label{ch:PhyXMiniProblems_problem_phyx_mini_0715}

\paragraph{Problem source.}
\#\# Physical scenario

A less-than-successful inventor wants to launch small satellites into orbit by launching them straight up from the surface of the earth at very high speed.

\#\# Question

By what percentage would your answer be in error if you used a flat-earth approximation?

\#\# Answer choices

- A: A: 2.1\%

- B: B: 2.3\%

- C: C: 2.5\%

- D: D: 2.7\%

\#\# Figure caption (auxiliary; use the image as primary evidence)

The image is a diagram illustrating the movement of an object away from Earth along the vertical axis labeled \textbackslash{}( y \textbackslash{}).

- At the bottom, there is a depiction of Earth with a radius labeled \textbackslash{}( R\_e \textbackslash{}).
- Below the Earth, there are two labeled positions: "Before" and "After."
- At "Before," the object (\textbackslash{}( y\_1 \textbackslash{})) is at \textbackslash{}( 0 \textbackslash{}) km, with velocity \textbackslash{}( v\_1 \textbackslash{}).
- "After" shows the object (\textbackslash{}( y\_2 \textbackslash{})) at \textbackslash{}( 400 \textbackslash{}) km, with velocity \textbackslash{}( v\_2 = 500 \textbackslash{}) m/s.
- Both positions show upward-pointing green arrows, indicating the direction of motion.
- Text labels: 
  - \textbackslash{}( y\_1 = 0 \textbackslash{}text\{ km\} \textbackslash{})
  - \textbackslash{}( y\_2 = 400 \textbackslash{}text\{ km\} \textbackslash{})
  - \textbackslash{}( v\_2 = 500 \textbackslash{}text\{ m/s\} \textbackslash{})

\#\# Dataset metadata

Dynamics; Physical Model Grounding Reasoning, Implicit Condition Reasoning

\paragraph{Recorded answer/context.}
C: C: 2.5\%

\paragraph{Figure/image path.}
/root/proposal\_for\_physic/science-mango/phyx\_mini\_run/phyx\_data/test\_image/715.png

\paragraph{Formalization target.}
create a compiling Lean file with sorry bodies at `PhyXMiniProblems/problem\_phyx\_mini\_0715.lean`. The Lean declarations must preserve the physical quantities, dimensions or dimensional roles, figure labels, governing-law hypotheses, and final relation expressed by this problem.
Use Mathlib/Physlib names found through LeanExplore where available. If a physics API is missing, introduce faithful local abstractions rather than scalar placeholder aliases.

\begin{theorem}[Physics formalization target]
\label{thm:physics:phyx_mini_0715:target}
\lean{PhyXMiniProblems.ProblemPhyXMini0715.flatEarthApproximationErrorIsApproximatelyThreePercentAndChoiceD}
\uses{def:physics:phyx-mini-0715:phyxminiproblems-problemphyxmini0715-speedinmeterspersecond, def:physics:phyx-mini-0715:phyxminiproblems-problemphyxmini0715-verticalsatellitelaunchsetup, def:physics:phyx-mini-0715:phyxminiproblems-problemphyxmini0715-matchesstraightupsatellitescenario, def:physics:phyx-mini-0715:phyxminiproblems-problemphyxmini0715-matchesdisplayedlaunchreadouts, def:physics:phyx-mini-0715:phyxminiproblems-problemphyxmini0715-matchesprimaryearthlaunchfigure, def:physics:phyx-mini-0715:phyxminiproblems-problemphyxmini0715-usesstandardearthcalibration, def:physics:phyx-mini-0715:phyxminiproblems-problemphyxmini0715-usessameboundarydataforflatapproximation, def:physics:phyx-mini-0715:phyxminiproblems-problemphyxmini0715-hasphysicalverticallaunchparameters, def:physics:phyx-mini-0715:phyxminiproblems-problemphyxmini0715-satisfiessatellitekineticenergylaw, def:physics:phyx-mini-0715:phyxminiproblems-problemphyxmini0715-satisfiessphericalearthpotentiallaw, def:physics:phyx-mini-0715:phyxminiproblems-problemphyxmini0715-satisfiesflatearthpotentiallaw, def:physics:phyx-mini-0715:phyxminiproblems-problemphyxmini0715-satisfiesmechanicalenergyconservation, def:physics:phyx-mini-0715:phyxminiproblems-problemphyxmini0715-sphericalpredictedinitialspeedinmeterspersecond, def:physics:phyx-mini-0715:phyxminiproblems-problemphyxmini0715-flatpredictedinitialspeedinmeterspersecond, def:physics:phyx-mini-0715:phyxminiproblems-problemphyxmini0715-flatearthpercentageerror, def:physics:phyx-mini-0715:phyxminiproblems-problemphyxmini0715-isclosestpercentagechoice}
The assigned autoformalize agent should translate this physics problem into Lean declarations in the covered file, with theorem and lemma proof bodies written as `by sorry`.
\end{theorem}
\begin{proof}
This is an autoformalization task, not a proof task. Produce statements that can later be proved by the physics prover without weakening the physical model.
\end{proof}

% --- Archon declaration topology begin ---
\section{Declaration topology}
\begin{definition}[Length Quantity]
  \label{def:physics:phyx-mini-0715:phyxminiproblems-problemphyxmini0715-lengthquantity}
  \lean{PhyXMiniProblems.ProblemPhyXMini0715.LengthQuantity}
  A nonnegative, unit-independent physical length
\end{definition}
\begin{proof}
  This is the declared model definition of Length Quantity.
\end{proof}
\begin{definition}[Mass Quantity]
  \label{def:physics:phyx-mini-0715:phyxminiproblems-problemphyxmini0715-massquantity}
  \lean{PhyXMiniProblems.ProblemPhyXMini0715.MassQuantity}
  A nonnegative, unit-independent physical mass
\end{definition}
\begin{proof}
  This is the declared model definition of Mass Quantity.
\end{proof}
\begin{definition}[Acceleration Quantity]
  \label{def:physics:phyx-mini-0715:phyxminiproblems-problemphyxmini0715-accelerationquantity}
  \lean{PhyXMiniProblems.ProblemPhyXMini0715.AccelerationQuantity}
  A nonnegative, unit-independent acceleration magnitude
\end{definition}
\begin{proof}
  This is the declared model definition of Acceleration Quantity.
\end{proof}
\begin{definition}[length Readout]
  \label{def:physics:phyx-mini-0715:phyxminiproblems-problemphyxmini0715-lengthreadout}
  \lean{PhyXMiniProblems.ProblemPhyXMini0715.lengthReadout}
  \uses{def:physics:phyx-mini-0715:phyxminiproblems-problemphyxmini0715-lengthquantity}
  Read a physical length in a selected length unit
\end{definition}
\begin{proof}
  This is the declared model definition of length Readout.
\end{proof}
\begin{definition}[length In Meters]
  \label{def:physics:phyx-mini-0715:phyxminiproblems-problemphyxmini0715-lengthinmeters}
  \lean{PhyXMiniProblems.ProblemPhyXMini0715.lengthInMeters}
  \uses{def:physics:phyx-mini-0715:phyxminiproblems-problemphyxmini0715-lengthquantity, def:physics:phyx-mini-0715:phyxminiproblems-problemphyxmini0715-lengthreadout}
  Read a physical length in metres
\end{definition}
\begin{proof}
  This is the declared model definition of length In Meters.
\end{proof}
\begin{definition}[length In Kilometers]
  \label{def:physics:phyx-mini-0715:phyxminiproblems-problemphyxmini0715-lengthinkilometers}
  \lean{PhyXMiniProblems.ProblemPhyXMini0715.lengthInKilometers}
  \uses{def:physics:phyx-mini-0715:phyxminiproblems-problemphyxmini0715-lengthquantity, def:physics:phyx-mini-0715:phyxminiproblems-problemphyxmini0715-lengthreadout}
  Read a physical length in kilometres
\end{definition}
\begin{proof}
  This is the declared model definition of length In Kilometers.
\end{proof}
\begin{definition}[mass In Kilograms]
  \label{def:physics:phyx-mini-0715:phyxminiproblems-problemphyxmini0715-massinkilograms}
  \lean{PhyXMiniProblems.ProblemPhyXMini0715.massInKilograms}
  \uses{def:physics:phyx-mini-0715:phyxminiproblems-problemphyxmini0715-massquantity}
  Read a physical mass in kilograms
\end{definition}
\begin{proof}
  This is the declared model definition of mass In Kilograms.
\end{proof}
\begin{definition}[speed In Meters Per Second]
  \label{def:physics:phyx-mini-0715:phyxminiproblems-problemphyxmini0715-speedinmeterspersecond}
  \lean{PhyXMiniProblems.ProblemPhyXMini0715.speedInMetersPerSecond}
  Read a speed in metres per second
\end{definition}
\begin{proof}
  This is the declared model definition of speed In Meters Per Second.
\end{proof}
\begin{definition}[acceleration In Meters Per Second Squared]
  \label{def:physics:phyx-mini-0715:phyxminiproblems-problemphyxmini0715-accelerationinmeterspersecondsquared}
  \lean{PhyXMiniProblems.ProblemPhyXMini0715.accelerationInMetersPerSecondSquared}
  \uses{def:physics:phyx-mini-0715:phyxminiproblems-problemphyxmini0715-accelerationquantity}
  Read an acceleration in metres per second squared
\end{definition}
\begin{proof}
  This is the declared model definition of acceleration In Meters Per Second Squared.
\end{proof}
\begin{definition}[energy In Joules]
  \label{def:physics:phyx-mini-0715:phyxminiproblems-problemphyxmini0715-energyinjoules}
  \lean{PhyXMiniProblems.ProblemPhyXMini0715.energyInJoules}
  Read an energy in joules
\end{definition}
\begin{proof}
  This is the declared model definition of energy In Joules.
\end{proof}
\begin{definition}[Launch Instant]
  \label{def:physics:phyx-mini-0715:phyxminiproblems-problemphyxmini0715-launchinstant}
  \lean{PhyXMiniProblems.ProblemPhyXMini0715.LaunchInstant}
  The two labelled states in the launch diagram
\end{definition}
\begin{proof}
  This is the declared model definition of Launch Instant.
\end{proof}
\begin{definition}[Gravity Model]
  \label{def:physics:phyx-mini-0715:phyxminiproblems-problemphyxmini0715-gravitymodel}
  \lean{PhyXMiniProblems.ProblemPhyXMini0715.GravityModel}
  The two gravitational models whose launch-speed predictions are compared
\end{definition}
\begin{proof}
  This is the declared model definition of Gravity Model.
\end{proof}
\begin{definition}[Launched Object Kind]
  \label{def:physics:phyx-mini-0715:phyxminiproblems-problemphyxmini0715-launchedobjectkind}
  \lean{PhyXMiniProblems.ProblemPhyXMini0715.LaunchedObjectKind}
  The kind of launched object described by the problem
\end{definition}
\begin{proof}
  This is the declared model definition of Launched Object Kind.
\end{proof}
\begin{definition}[Central Body]
  \label{def:physics:phyx-mini-0715:phyxminiproblems-problemphyxmini0715-centralbody}
  \lean{PhyXMiniProblems.ProblemPhyXMini0715.CentralBody}
  The central body pictured below the launch path
\end{definition}
\begin{proof}
  This is the declared model definition of Central Body.
\end{proof}
\begin{definition}[Launch Path]
  \label{def:physics:phyx-mini-0715:phyxminiproblems-problemphyxmini0715-launchpath}
  \lean{PhyXMiniProblems.ProblemPhyXMini0715.LaunchPath}
  The path used by the inventor to launch the satellite
\end{definition}
\begin{proof}
  This is the declared model definition of Launch Path.
\end{proof}
\begin{definition}[Vertical Direction]
  \label{def:physics:phyx-mini-0715:phyxminiproblems-problemphyxmini0715-verticaldirection}
  \lean{PhyXMiniProblems.ProblemPhyXMini0715.VerticalDirection}
  Directions available along the vertical launch axis
\end{definition}
\begin{proof}
  This is the declared model definition of Vertical Direction.
\end{proof}
\begin{definition}[Axis Symbol]
  \label{def:physics:phyx-mini-0715:phyxminiproblems-problemphyxmini0715-axissymbol}
  \lean{PhyXMiniProblems.ProblemPhyXMini0715.AxisSymbol}
  Coordinate-axis symbols that may be printed in the figure
\end{definition}
\begin{proof}
  This is the declared model definition of Axis Symbol.
\end{proof}
\begin{definition}[Velocity Symbol]
  \label{def:physics:phyx-mini-0715:phyxminiproblems-problemphyxmini0715-velocitysymbol}
  \lean{PhyXMiniProblems.ProblemPhyXMini0715.VelocitySymbol}
  Velocity labels attached to the two satellite markers
\end{definition}
\begin{proof}
  This is the declared model definition of Velocity Symbol.
\end{proof}
\begin{definition}[Altitude Symbol]
  \label{def:physics:phyx-mini-0715:phyxminiproblems-problemphyxmini0715-altitudesymbol}
  \lean{PhyXMiniProblems.ProblemPhyXMini0715.AltitudeSymbol}
  Altitude-coordinate labels attached to the two satellite markers
\end{definition}
\begin{proof}
  This is the declared model definition of Altitude Symbol.
\end{proof}
\begin{definition}[Radius Symbol]
  \label{def:physics:phyx-mini-0715:phyxminiproblems-problemphyxmini0715-radiussymbol}
  \lean{PhyXMiniProblems.ProblemPhyXMini0715.RadiusSymbol}
  The Earth-radius symbol printed inside the central body
\end{definition}
\begin{proof}
  This is the declared model definition of Radius Symbol.
\end{proof}
\begin{definition}[Earth Launch Figure]
  \label{def:physics:phyx-mini-0715:phyxminiproblems-problemphyxmini0715-earthlaunchfigure}
  \lean{PhyXMiniProblems.ProblemPhyXMini0715.EarthLaunchFigure}
  \uses{def:physics:phyx-mini-0715:phyxminiproblems-problemphyxmini0715-launchinstant, def:physics:phyx-mini-0715:phyxminiproblems-problemphyxmini0715-centralbody, def:physics:phyx-mini-0715:phyxminiproblems-problemphyxmini0715-axissymbol, def:physics:phyx-mini-0715:phyxminiproblems-problemphyxmini0715-velocitysymbol, def:physics:phyx-mini-0715:phyxminiproblems-problemphyxmini0715-altitudesymbol, def:physics:phyx-mini-0715:phyxminiproblems-problemphyxmini0715-radiussymbol}
  Literal evidence visible in 715.png. The scalar fields are printed readouts and are kept separate from the physical dimensionful quantities
\end{definition}
\begin{proof}
  This is the declared model definition of Earth Launch Figure.
\end{proof}
\begin{definition}[Vertical Satellite Launch Setup]
  \label{def:physics:phyx-mini-0715:phyxminiproblems-problemphyxmini0715-verticalsatellitelaunchsetup}
  \lean{PhyXMiniProblems.ProblemPhyXMini0715.VerticalSatelliteLaunchSetup}
  \uses{def:physics:phyx-mini-0715:phyxminiproblems-problemphyxmini0715-lengthquantity, def:physics:phyx-mini-0715:phyxminiproblems-problemphyxmini0715-massquantity, def:physics:phyx-mini-0715:phyxminiproblems-problemphyxmini0715-accelerationquantity, def:physics:phyx-mini-0715:phyxminiproblems-problemphyxmini0715-launchinstant, def:physics:phyx-mini-0715:phyxminiproblems-problemphyxmini0715-gravitymodel, def:physics:phyx-mini-0715:phyxminiproblems-problemphyxmini0715-launchedobjectkind, def:physics:phyx-mini-0715:phyxminiproblems-problemphyxmini0715-centralbody, def:physics:phyx-mini-0715:phyxminiproblems-problemphyxmini0715-launchpath, def:physics:phyx-mini-0715:phyxminiproblems-problemphyxmini0715-verticaldirection, def:physics:phyx-mini-0715:phyxminiproblems-problemphyxmini0715-earthlaunchfigure}
  Independent physical quantities for the true spherical model and the flat-Earth approximation. In particular, both initial speeds are independent fields; neither is defined from an answer choice or predicted formula
\end{definition}
\begin{proof}
  This is the declared model definition of Vertical Satellite Launch Setup.
\end{proof}
\begin{definition}[Matches Straight Up Satellite Scenario]
  \label{def:physics:phyx-mini-0715:phyxminiproblems-problemphyxmini0715-matchesstraightupsatellitescenario}
  \lean{PhyXMiniProblems.ProblemPhyXMini0715.MatchesStraightUpSatelliteScenario}
  \uses{def:physics:phyx-mini-0715:phyxminiproblems-problemphyxmini0715-verticalsatellitelaunchsetup}
  Qualitative facts stated in the problem prose
\end{definition}
\begin{proof}
  This is the declared model definition of Matches Straight Up Satellite Scenario.
\end{proof}
\begin{definition}[Matches Displayed Launch Readouts]
  \label{def:physics:phyx-mini-0715:phyxminiproblems-problemphyxmini0715-matchesdisplayedlaunchreadouts}
  \lean{PhyXMiniProblems.ProblemPhyXMini0715.MatchesDisplayedLaunchReadouts}
  \uses{def:physics:phyx-mini-0715:phyxminiproblems-problemphyxmini0715-lengthinkilometers, def:physics:phyx-mini-0715:phyxminiproblems-problemphyxmini0715-speedinmeterspersecond, def:physics:phyx-mini-0715:phyxminiproblems-problemphyxmini0715-verticalsatellitelaunchsetup}
  The numerical state data printed in the primary image
\end{definition}
\begin{proof}
  This is the declared model definition of Matches Displayed Launch Readouts.
\end{proof}
\begin{definition}[Matches Primary Earth Launch Figure]
  \label{def:physics:phyx-mini-0715:phyxminiproblems-problemphyxmini0715-matchesprimaryearthlaunchfigure}
  \lean{PhyXMiniProblems.ProblemPhyXMini0715.MatchesPrimaryEarthLaunchFigure}
  \uses{def:physics:phyx-mini-0715:phyxminiproblems-problemphyxmini0715-lengthinkilometers, def:physics:phyx-mini-0715:phyxminiproblems-problemphyxmini0715-speedinmeterspersecond, def:physics:phyx-mini-0715:phyxminiproblems-problemphyxmini0715-verticalsatellitelaunchsetup}
  Raster-visible geometry and labels. The before-state velocity has the symbol v₁ but no numerical value, whereas the after-state label is v₂ = 500 m/s
\end{definition}
\begin{proof}
  This is the declared model definition of Matches Primary Earth Launch Figure.
\end{proof}
\begin{definition}[Uses Standard Earth Calibration]
  \label{def:physics:phyx-mini-0715:phyxminiproblems-problemphyxmini0715-usesstandardearthcalibration}
  \lean{PhyXMiniProblems.ProblemPhyXMini0715.UsesStandardEarthCalibration}
  \uses{def:physics:phyx-mini-0715:phyxminiproblems-problemphyxmini0715-lengthinkilometers, def:physics:phyx-mini-0715:phyxminiproblems-problemphyxmini0715-accelerationinmeterspersecondsquared, def:physics:phyx-mini-0715:phyxminiproblems-problemphyxmini0715-verticalsatellitelaunchsetup}
  Standard numerical data needed because the raster labels Rₑ but does not print its value. The mean Earth radius and conventional standard gravity are kept separate from facts directly read from the source
\end{definition}
\begin{proof}
  This is the declared model definition of Uses Standard Earth Calibration.
\end{proof}
\begin{definition}[Uses Same Boundary Data For Flat Approximation]
  \label{def:physics:phyx-mini-0715:phyxminiproblems-problemphyxmini0715-usessameboundarydataforflatapproximation}
  \lean{PhyXMiniProblems.ProblemPhyXMini0715.UsesSameBoundaryDataForFlatApproximation}
  \uses{def:physics:phyx-mini-0715:phyxminiproblems-problemphyxmini0715-speedinmeterspersecond, def:physics:phyx-mini-0715:phyxminiproblems-problemphyxmini0715-verticalsatellitelaunchsetup}
  The flat model answers the same boundary-value question as the spherical model: it uses the same final altitude and final speed. Altitudes are already shared fields, so only the independently stored final speeds must be related
\end{definition}
\begin{proof}
  This is the declared model definition of Uses Same Boundary Data For Flat Approximation.
\end{proof}
\begin{definition}[Has Physical Vertical Launch Parameters]
  \label{def:physics:phyx-mini-0715:phyxminiproblems-problemphyxmini0715-hasphysicalverticallaunchparameters}
  \lean{PhyXMiniProblems.ProblemPhyXMini0715.HasPhysicalVerticalLaunchParameters}
  \uses{def:physics:phyx-mini-0715:phyxminiproblems-problemphyxmini0715-lengthinmeters, def:physics:phyx-mini-0715:phyxminiproblems-problemphyxmini0715-massinkilograms, def:physics:phyx-mini-0715:phyxminiproblems-problemphyxmini0715-speedinmeterspersecond, def:physics:phyx-mini-0715:phyxminiproblems-problemphyxmini0715-accelerationinmeterspersecondsquared, def:physics:phyx-mini-0715:phyxminiproblems-problemphyxmini0715-verticalsatellitelaunchsetup}
  Positivity and nondegeneracy conditions for the physical comparison
\end{definition}
\begin{proof}
  This is the declared model definition of Has Physical Vertical Launch Parameters.
\end{proof}
\begin{definition}[Satisfies Satellite Kinetic Energy Law]
  \label{def:physics:phyx-mini-0715:phyxminiproblems-problemphyxmini0715-satisfiessatellitekineticenergylaw}
  \lean{PhyXMiniProblems.ProblemPhyXMini0715.SatisfiesSatelliteKineticEnergyLaw}
  \uses{def:physics:phyx-mini-0715:phyxminiproblems-problemphyxmini0715-massinkilograms, def:physics:phyx-mini-0715:phyxminiproblems-problemphyxmini0715-speedinmeterspersecond, def:physics:phyx-mini-0715:phyxminiproblems-problemphyxmini0715-energyinjoules, def:physics:phyx-mini-0715:phyxminiproblems-problemphyxmini0715-verticalsatellitelaunchsetup}
  Translational kinetic energy K = (1/2) m v² in either model
\end{definition}
\begin{proof}
  This is the declared model definition of Satisfies Satellite Kinetic Energy Law.
\end{proof}
\begin{definition}[Satisfies Spherical Earth Potential Law]
  \label{def:physics:phyx-mini-0715:phyxminiproblems-problemphyxmini0715-satisfiessphericalearthpotentiallaw}
  \lean{PhyXMiniProblems.ProblemPhyXMini0715.SatisfiesSphericalEarthPotentialLaw}
  \uses{def:physics:phyx-mini-0715:phyxminiproblems-problemphyxmini0715-lengthinmeters, def:physics:phyx-mini-0715:phyxminiproblems-problemphyxmini0715-massinkilograms, def:physics:phyx-mini-0715:phyxminiproblems-problemphyxmini0715-accelerationinmeterspersecondsquared, def:physics:phyx-mini-0715:phyxminiproblems-problemphyxmini0715-energyinjoules, def:physics:phyx-mini-0715:phyxminiproblems-problemphyxmini0715-verticalsatellitelaunchsetup}
  The inverse-square Newtonian potential, normalized to zero at infinity and written using surface gravity g₀ = GM / Rₑ²: U(y) = -m g₀ Rₑ² / (Rₑ + y)
\end{definition}
\begin{proof}
  This is the declared model definition of Satisfies Spherical Earth Potential Law.
\end{proof}
\begin{definition}[Satisfies Flat Earth Potential Law]
  \label{def:physics:phyx-mini-0715:phyxminiproblems-problemphyxmini0715-satisfiesflatearthpotentiallaw}
  \lean{PhyXMiniProblems.ProblemPhyXMini0715.SatisfiesFlatEarthPotentialLaw}
  \uses{def:physics:phyx-mini-0715:phyxminiproblems-problemphyxmini0715-lengthinmeters, def:physics:phyx-mini-0715:phyxminiproblems-problemphyxmini0715-massinkilograms, def:physics:phyx-mini-0715:phyxminiproblems-problemphyxmini0715-accelerationinmeterspersecondsquared, def:physics:phyx-mini-0715:phyxminiproblems-problemphyxmini0715-energyinjoules, def:physics:phyx-mini-0715:phyxminiproblems-problemphyxmini0715-verticalsatellitelaunchsetup}
  The flat-Earth approximation uses a uniform downward acceleration and the near-surface potential U\_flat(y) = m g₀ y
\end{definition}
\begin{proof}
  This is the declared model definition of Satisfies Flat Earth Potential Law.
\end{proof}
\begin{definition}[total Mechanical Energy In Joules]
  \label{def:physics:phyx-mini-0715:phyxminiproblems-problemphyxmini0715-totalmechanicalenergyinjoules}
  \lean{PhyXMiniProblems.ProblemPhyXMini0715.totalMechanicalEnergyInJoules}
  \uses{def:physics:phyx-mini-0715:phyxminiproblems-problemphyxmini0715-energyinjoules, def:physics:phyx-mini-0715:phyxminiproblems-problemphyxmini0715-launchinstant, def:physics:phyx-mini-0715:phyxminiproblems-problemphyxmini0715-gravitymodel, def:physics:phyx-mini-0715:phyxminiproblems-problemphyxmini0715-verticalsatellitelaunchsetup}
  Total mechanical energy in one model at one labelled instant
\end{definition}
\begin{proof}
  This is the declared model definition of total Mechanical Energy In Joules.
\end{proof}
\begin{definition}[Satisfies Mechanical Energy Conservation]
  \label{def:physics:phyx-mini-0715:phyxminiproblems-problemphyxmini0715-satisfiesmechanicalenergyconservation}
  \lean{PhyXMiniProblems.ProblemPhyXMini0715.SatisfiesMechanicalEnergyConservation}
  \uses{def:physics:phyx-mini-0715:phyxminiproblems-problemphyxmini0715-verticalsatellitelaunchsetup, def:physics:phyx-mini-0715:phyxminiproblems-problemphyxmini0715-totalmechanicalenergyinjoules}
  Each conservative model equates the before and after mechanical energies
\end{definition}
\begin{proof}
  This is the declared model definition of Satisfies Mechanical Energy Conservation.
\end{proof}
\begin{definition}[spherical Predicted Initial Speed In Meters Per Second]
  \label{def:physics:phyx-mini-0715:phyxminiproblems-problemphyxmini0715-sphericalpredictedinitialspeedinmeterspersecond}
  \lean{PhyXMiniProblems.ProblemPhyXMini0715.sphericalPredictedInitialSpeedInMetersPerSecond}
  \uses{def:physics:phyx-mini-0715:phyxminiproblems-problemphyxmini0715-lengthinmeters, def:physics:phyx-mini-0715:phyxminiproblems-problemphyxmini0715-speedinmeterspersecond, def:physics:phyx-mini-0715:phyxminiproblems-problemphyxmini0715-accelerationinmeterspersecondsquared, def:physics:phyx-mini-0715:phyxminiproblems-problemphyxmini0715-verticalsatellitelaunchsetup}
  Launch-speed readout predicted by the spherical inverse-square model
\end{definition}
\begin{proof}
  This is the declared model definition of spherical Predicted Initial Speed In Meters Per Second.
\end{proof}
\begin{definition}[flat Predicted Initial Speed In Meters Per Second]
  \label{def:physics:phyx-mini-0715:phyxminiproblems-problemphyxmini0715-flatpredictedinitialspeedinmeterspersecond}
  \lean{PhyXMiniProblems.ProblemPhyXMini0715.flatPredictedInitialSpeedInMetersPerSecond}
  \uses{def:physics:phyx-mini-0715:phyxminiproblems-problemphyxmini0715-lengthinmeters, def:physics:phyx-mini-0715:phyxminiproblems-problemphyxmini0715-speedinmeterspersecond, def:physics:phyx-mini-0715:phyxminiproblems-problemphyxmini0715-accelerationinmeterspersecondsquared, def:physics:phyx-mini-0715:phyxminiproblems-problemphyxmini0715-verticalsatellitelaunchsetup}
  Launch-speed readout predicted by the constant-g flat-Earth model
\end{definition}
\begin{proof}
  This is the declared model definition of flat Predicted Initial Speed In Meters Per Second.
\end{proof}
\begin{definition}[flat Earth Percentage Error]
  \label{def:physics:phyx-mini-0715:phyxminiproblems-problemphyxmini0715-flatearthpercentageerror}
  \lean{PhyXMiniProblems.ProblemPhyXMini0715.flatEarthPercentageError}
  \uses{def:physics:phyx-mini-0715:phyxminiproblems-problemphyxmini0715-speedinmeterspersecond, def:physics:phyx-mini-0715:phyxminiproblems-problemphyxmini0715-verticalsatellitelaunchsetup}
  Standard absolute percentage error of the approximate launch speed, using the spherical-model launch speed as the exact reference value
\end{definition}
\begin{proof}
  This is the declared model definition of flat Earth Percentage Error.
\end{proof}
\begin{definition}[Answer Choice]
  \label{def:physics:phyx-mini-0715:phyxminiproblems-problemphyxmini0715-answerchoice}
  \lean{PhyXMiniProblems.ProblemPhyXMini0715.AnswerChoice}
  The four percentage labels displayed by the problem
\end{definition}
\begin{proof}
  This is the declared model definition of Answer Choice.
\end{proof}
\begin{definition}[percentage]
  \label{def:physics:phyx-mini-0715:phyxminiproblems-problemphyxmini0715-answerchoice-percentage}
  \lean{PhyXMiniProblems.ProblemPhyXMini0715.AnswerChoice.percentage}
  \uses{def:physics:phyx-mini-0715:phyxminiproblems-problemphyxmini0715-answerchoice}
  Percentage printed next to an answer label
\end{definition}
\begin{proof}
  This is the declared model definition of percentage.
\end{proof}
\begin{definition}[recorded Dataset Answer]
  \label{def:physics:phyx-mini-0715:phyxminiproblems-problemphyxmini0715-recordeddatasetanswer}
  \lean{PhyXMiniProblems.ProblemPhyXMini0715.recordedDatasetAnswer}
  \uses{def:physics:phyx-mini-0715:phyxminiproblems-problemphyxmini0715-answerchoice}
  Dataset answer label retained as metadata, never used as a premise
\end{definition}
\begin{proof}
  This is the declared model definition of recorded Dataset Answer.
\end{proof}
\begin{definition}[Is Closest Percentage Choice]
  \label{def:physics:phyx-mini-0715:phyxminiproblems-problemphyxmini0715-isclosestpercentagechoice}
  \lean{PhyXMiniProblems.ProblemPhyXMini0715.IsClosestPercentageChoice}
  \uses{def:physics:phyx-mini-0715:phyxminiproblems-problemphyxmini0715-answerchoice}
  A displayed choice is at least as close as every option to the model error
\end{definition}
\begin{proof}
  This is the declared model definition of Is Closest Percentage Choice.
\end{proof}
\begin{lemma}[initial Speeds eq model Predictions]
  \label{lem:physics:phyx-mini-0715:phyxminiproblems-problemphyxmini0715-initialspeeds-eq-modelpredictions}
  \lean{PhyXMiniProblems.ProblemPhyXMini0715.initialSpeeds_eq_modelPredictions}
  \uses{def:physics:phyx-mini-0715:phyxminiproblems-problemphyxmini0715-speedinmeterspersecond, def:physics:phyx-mini-0715:phyxminiproblems-problemphyxmini0715-verticalsatellitelaunchsetup, def:physics:phyx-mini-0715:phyxminiproblems-problemphyxmini0715-hasphysicalverticallaunchparameters, def:physics:phyx-mini-0715:phyxminiproblems-problemphyxmini0715-satisfiessatellitekineticenergylaw, def:physics:phyx-mini-0715:phyxminiproblems-problemphyxmini0715-satisfiessphericalearthpotentiallaw, def:physics:phyx-mini-0715:phyxminiproblems-problemphyxmini0715-satisfiesflatearthpotentiallaw, def:physics:phyx-mini-0715:phyxminiproblems-problemphyxmini0715-satisfiesmechanicalenergyconservation, def:physics:phyx-mini-0715:phyxminiproblems-problemphyxmini0715-sphericalpredictedinitialspeedinmeterspersecond, def:physics:phyx-mini-0715:phyxminiproblems-problemphyxmini0715-flatpredictedinitialspeedinmeterspersecond}
  Energy conservation and the two potential laws give the respective spherical and flat-Earth launch-speed formulas. This is a symbolic derived conclusion; it contains no answer-choice percentage
\end{lemma}
\begin{proof}
  The conclusion follows from the governing relations and figure readouts named in the statement of initial Speeds eq model Predictions.
\end{proof}
% --- Archon declaration topology end ---
% --- Archon physics formalization source end ---
